\documentclass[nocfonts]{nostarch}

\usepackage{xcolor}

\usepackage{csquotes}
\usepackage[style=verbose-note]{biblatex}
\usepackage{subfiles}

% IMPORTANT: babel before tikz/pgfplots
\usepackage[spanish, es-noshorthands]{babel}
\usepackage{tikz}
\usetikzlibrary{calc,positioning, arrows.meta}
\usetikzlibrary{shapes, arrows}
\usetikzlibrary{backgrounds}
\tikzset{
    block/.style = {
        draw,
        rectangle,
        minimum width=3cm,
        minimum height=1.5cm,
        align=center
    }
}
\usetikzlibrary{shapes.geometric, arrows}

\tikzstyle{startstop} = [rectangle, rounded corners, 
minimum width=3cm, 
minimum height=1cm,
text centered, 
draw=black, 
fill=red!30]

\tikzstyle{io} = [trapezium, 
trapezium stretches=true, % A later addition
trapezium left angle=70, 
trapezium right angle=110, 
minimum width=3cm, 
minimum height=1cm, text centered, 
draw=black, fill=blue!30]

\tikzstyle{process} = [rectangle, 
minimum width=3cm, 
minimum height=1cm, 
text centered, 
text width=3cm, 
draw=black, 
fill=orange!30]

\tikzstyle{decision} = [diamond, 
minimum width=3cm, 
minimum height=1cm, 
text centered, 
draw=black, 
fill=green!30]
\tikzstyle{arrow} = [thick,->,>=stealth]


\usepackage{circuitikz}
\usepackage{pgfplots}
\pgfplotsset{compat=1.18}
\usepackage{nshyper}

\begin{document}

\frontmatter

\title{Máquinas sonoras}

\subtitle{Electrónicas y computacionales}

\author{Aarón Montoya Moraga y Matías Ignacio Serrano Acevedo}

\nostarchlocation{Santiago}

\makehalftitle

\maketitle

\brieftableofcontents

\tableofcontents

\mainmatter

\subfile{info/prefacio}
\subfile{info/autores}
\subfile{info/cursos}
\subfile{info/herramientas}
\subfile{info/agradecimientos}

\subfile{maquinas/energias}
\subfile{maquinas/relojes}
\subfile{maquinas/secuenciadores}



\backmatter



% \part{Conceptos maquinales}

% \subfile{maquinas/osciladores}
% \subfile{maquinas/amplificadores}
% \subfile{maquinas/moduladores}
% \subfile{maquinas/filtros}


% \part{Electricidad}

% \subfile{electricidad/voltaje}
% \subfile{electricidad/corriente}
% \subfile{electricidad/potencia}
% \subfile{electricidad/fuentes-de-poder.tex}

% \part{Electrónica}

% \chapter{Componentes básicos}

% \subfile{componentes-basicos/placa-de-pruebas}
% \subfile{componentes-basicos/resistores}
% \subfile{componentes-basicos/capacitores}

% \chapter{Conectores}

% \chapter{Circuitos integrados}

% \subfile{circuitos-integrados/chip555}
% \subfile{circuitos-integrados/chip4017}
% \subfile{circuitos-integrados/chip4093}

% \part{Computación}

% \chapter{Bloques fundamentales de la programación}

% \subfile{computacion/variables}
% \subfile{computacion/funciones}

% \part{Interfaces}

% \subfile{interfaces/botones}


% \part{Ejemplos}

% \subfile{ejemplos/cuatro-cuadradas}

% \part{Anexos}

% \subfile{anexos/albumes}
% \subfile{anexos/maquinas}
% \subfile{anexos/organizaciones}
% \subfile{anexos/personas}


% \printglossaries

% \printbibliography

\end{document}
